\section{Related Work}
\label{sec:related_work}

Scene text replacement and video text editing have been studied across multiple subfields, including image-based scene text editing, video generation, and optical flow-based tracking.

\subsection{Image-Based Scene Text Editing}

Early work in scene text editing primarily focuses on static images. Methods such as CLASTE \cite{claste} demonstrate strong performance in replacing text while preserving font style, texture, and local appearance consistency. More recent approaches, such as TextMastero \cite{textmastero}, further improve multilingual support and visual quality through advanced generative modeling techniques.

In addition, font transfer methods such as Few-Shot Font Style Transfer \cite{fewshotfont} and Fontify \cite{fontify} address the challenge of style consistency across languages, enabling more realistic cross-lingual text rendering.

These methods operate in single-image settings and do not model temporal coherence, which limits their applicability to video scenarios.

\subsection{Video Generation and Editing Models}

Recent large-scale generative video models have demonstrated strong capabilities in producing visually coherent and high-quality video content. These models excel in global temporal consistency and semantic alignment across frames.

However, they lack explicit control over fine-grained textual structure. As a result, text rendered in generated videos is often distorted or semantically incorrect, limiting their use for precise scene text replacement tasks.

\subsection{Tracking and Optical Flow-Based Methods}

Classical video processing pipelines rely on optical flow and feature tracking methods such as Lucas-Kanade, as well as modern learned trackers. These methods provide robust temporal correspondences and are widely used for propagating geometric information across frames.

However, they are not designed for semantic editing or appearance transformation tasks, and therefore must be combined with generative or editing models in practical systems.

\subsection{Video Scene Text Replacement Systems}

Existing end-to-end systems such as STRIVE \cite{strive} combine detection, frontalization, editing, and propagation in a multi-stage pipeline for video scene text replacement.

While such systems demonstrate the feasibility of structured pipelines for this task, they remain limited in handling challenging real-world conditions such as lighting variation, motion blur, and long-range temporal consistency.



Overall, prior work can be grouped into three directions: image-based text editing, video generation models, and tracking-based propagation methods. However, these approaches remain largely decoupled, motivating more integrated pipeline designs for robust cross-language scene text replacement in video.